\documentclass[10pt]{article}

\usepackage[utf8]{inputenc}

\usepackage[T1]{fontenc}

\usepackage{amsmath}

\usepackage{amsfonts}

\usepackage{amssymb}

\usepackage[version=4]{mhchem}

\usepackage{stmaryrd}

\usepackage{graphicx}

\usepackage[export]{adjustbox}

\graphicspath{ {./images/} }

\usepackage{mathrsfs}

\usepackage{bbold}



\begin{document}

дели, реализуемость по Клини, ступенчатая семантическая система А. А. Маркова.



Лит.: [1] К а р н а I Р., Значение и необходимость пер. с англ., М., 1959; [2] ч ё р ч А., Введение в математическую логику, пер. с англ., М., 1960; [3] К л и н и С. К., Введение в метаматематику, пер. с англ., М., 1957; [4] Д р а г а л и н А. Г., Математический интуитонизм. зательств, М., 1979; [5] Ф е й с Р., Модальная логика, пер. с СнгМИ.ЙНВАРИАНТ - 1) С.- то же, что полуинвариант. 2) С.- одна из числовых характеристик случайных величин, родственная понятию момента старшего порядка. Если $\xi=\left(\xi_{1}, \ldots, \xi_{k}\right)-$ случайный вектор, $\varphi_{\xi}(t)=\mathrm{E} e^{i(t, \xi)}-$ его характеристич. функция, $t=\left(t_{1}, \ldots, t_{k}\right), t_{i} \in R$,



\[

(t, \xi)=\sum_{i=1}^{k} t_{i} \xi_{i}

\]



и для нек-рого $n \geqslant 1$ моменты $\mathrm{E}\left|\xi_{i}\right|^{n}<\infty, i=1, \ldots, k$, то существуют (смешанные) моменты



\[

m_{\xi}^{\left(v_{1}, \ldots, v_{k}\right)}=\mathrm{E} \xi_{1}^{v_{1}} \ldots \xi v_{k}

\]



для всех неотрицательных целочисленных $v_{1}, \ldots, v_{k}$ таких, что $v_{1}+\ldots+v_{k} \leqslant n$. Тогда



\[

\begin{aligned}

& \varphi_{\xi}(t)=\sum_{v_{1}+\ldots+v_{k} \leqslant n} \frac{i^{v_{1}+\ldots+v_{k}}}{v_{1} ! \ldots v_{k} !} m_{\xi}^{\left(v_{1}, \ldots, v_{k}\right)} \times \\

& \times t_{1}^{v_{1}} \ldots t_{k}^{v_{k}}+o\left(|t|^{n}\right),

\end{aligned}

\]



где $|t|=\left|t_{1}\right|+\ldots+\left|t_{k}\right|$ и для достаточно малых $|t|$ главное значение $\ln \varphi_{\xi}(t)$ представимо по формуле Тейлора в виде



\[

\begin{aligned}

\ln \varphi_{\xi}(t)=\sum & \left.v_{1}+\ldots+v_{k} \leqslant n \frac{i^{v_{1}+\ldots+v_{k}}}{v_{1} ! \ldots v_{k} !} s_{\xi}^{\left(v_{1}\right.}, \ldots, v_{k}\right) \times \\

& \times t_{1}^{v_{1}} \ldots t_{k}^{v_{k}}+o\left(|t|^{n}\right),

\end{aligned}

\]



где коэффициенты $s_{\xi}^{\left(v_{1}, \ldots, v_{k}\right)}$ наз. (с м е ші а н н ы м и) с е ми и н. в а р и а н т а ми, или кум уля н т ам и, порядка $v=\left(v_{1}, \ldots, v_{k}\right)$ вектора $\xi=\left(\xi_{1}, \ldots, \xi_{k}\right)$. Для независимых случайных векторов $\xi=\left(\xi_{1}, \ldots\right.$



$\left.\ldots, \xi_{k}\right)$ и $\eta=\left(\eta_{1}, \ldots, \eta_{k}\right)$



\[

s_{\xi+\eta}^{\left(v_{1}, \ldots, v_{k}\right)}=s_{\xi}^{\left(v_{1}, \ldots, v_{k}\right)}+s_{\eta}^{\left(v_{1}, \ldots, v_{k}\right)},

\]



то есть С. суммы независимых случайных векторов есть сумма С. Именно это и послужило причиной термина "семиинвариант», отражающего свойство аддитивности для случая независимых величин (но әто свойство уже, вообще говоря, не верно для зависимых величин)



Между моментами и С. справедливы следующие форормулы связи:



\[

\begin{aligned}

& \left.m_{\xi}^{(v)}=\sum^{*} \frac{1}{q !} \frac{v !}{\lambda(1) ! \ldots \lambda(q) !} \prod_{p=1}^{q} s_{\xi}^{(\lambda(p)}\right),

\end{aligned}

\]



\begin{center}

\includegraphics[max width=\textwidth]{2023_07_08_caecfb7a26b7ee5d9ef1g-1}

\end{center}



где $\sum^{*}$ означает суммирование по всем ушорядоченным наборам целых неотрицательных векторов $\lambda^{(p)}$, $\left|\lambda^{(p)}\right|>0$, дающих в сумме вектор $v$. В частности, если $\xi$ - случайная величина $(k=1), m_{n}=m_{\xi}^{(n)}=\mathrm{E} \xi^{n}, s_{n}=$ $=s_{\xi}^{(n)}$, то



\[

\begin{aligned}

& m_{1}=s_{1}, \\

& m_{2}=s_{2}+s_{1}^{2}, \\

& m_{3}=s_{3}+3 s_{1} s_{2}+s_{1}^{3}, \\

& m_{4}=s_{4}+3 s_{2}^{2}+4 s_{1} s_{3}+6 s_{1}^{2} s_{2}+s_{1}^{4}, \\

& s_{1}=m_{1}(=\mathrm{E} \xi), \\

& s_{2}=m_{2}-m_{1}^{2}(=\mathrm{D} \xi), \\

& s_{3}=m_{3}-3 m_{1} m_{2}+2 m_{1}^{3}, \\

& s_{4}=m_{4}-3 m_{2}^{2}-4 m_{1} m_{3}+12 m_{1}^{2} m_{2}-6 m_{1}^{4} .

\end{aligned}

\]



Лum.: [1] Л е о н о в В. П., ШШ и р я е в А. Н., "Теория вероятн. и ее примен.", 1959, т. 4, в. 3, с. 342-55. А. Ширлев.



СЕМИМАРТИНГАЛ - стохастический процесс, представимый в виде суммы локального мартингала и процесса локально ограниченной вариации. При формальном определении С. исходят из допущения, что все рассмотрения ведутся на стохастич. базисе $(\Omega$, $\mathscr{F}, \mathbb{F}, \mathrm{P})$, где $\mathbb{F}=\left(\mathscr{F}_{t}\right)_{t \geqslant 0}$. Стохастич. процесс $X=$ $=\left(X_{t}, \mathscr{F}_{t}\right)_{t \geqslant 0}$ наз. с е ми м а р т и нг а л о м, если его траектории непрерывны справа и имеют пределы слева и он представим в виде $X_{t}=M_{t}+V_{t}$, где $M=$ $=\left(M_{t}, \mathscr{F}_{t}\right)$ - локальный мартингал, а $V=\left(V_{t}, \mathscr{F}_{t}\right)-$ процесс локально ограниченной вариации, т. е.



\[

\int_{0}^{t}\left|d V_{s}(\omega)\right|<\infty, t>0, \omega \in \Omega .

\]



Такое представление, вообще говоря, неоднозначно. Однако в классе разложений с предсказуемыми процессами $V$ рассматриваемое представление единственно (с точностью до стохастич. эквивалентности). К классу С. относятся (помимо локальных мартингалов и процессов с локально ограниченными вариациями) локальные супермартингалы и субмартингалы, процессы $X$ с независимыми приращениями, для к-рых функция $f(t)=\mathrm{E} e^{i \lambda X t}$ является функцией локально ограниченной вариации для любого $\lambda \in \mathbb{R}$ (и значит - все процессы со стационарными независимыми приращениями), процессы Ито, процессы диффузионного типа и др. Класс С. инвариантен относительно эквивалентной замены меры. Если $X$ есть С., а $f=f(x)$ - дважды непрерывно дифференцируемая функция, то $f(X)=$ $=\left(f\left(X_{t}\right), \mathscr{J}_{t}\right)$ также С. При этом (фо р м у л а И т о) $f\left(X_{t}\right)=f\left(X_{0}\right)+\int_{0}^{t} f^{\prime}\left(X_{s_{-}}\right) d X_{s}+\frac{1}{2} \int_{0}^{t} f^{\prime \prime}\left(X_{s_{-}}\right) d[X, X]_{s}^{c}+$



\[

+\sum_{0<s \leqslant t}\left[f\left(X_{s}\right)-f\left(X_{s^{-}}\right)-f^{\prime}\left(X_{s^{-}}\right) \Delta X_{s}\right]

\]



или, что эквивалентно,



$f\left(X_{t}\right)=f\left(X_{0}\right)+\int_{0}^{t} f^{\prime}\left(X_{s^{-}}\right) d X_{s}+\frac{1}{2} \int_{0}^{t} f^{\prime \prime}\left(X_{s^{-}}\right) d[X, X]_{s}+$



\[

+\sum_{0<s \leqslant t}\left[f\left(X_{s}\right)-f\left(X_{s^{-}}\right)-f^{\prime}\left(X_{s^{-}}\right) \Delta X_{s^{-}}\right.

\][



\textbackslash left.-\textbackslash frac\{1\}\{2\} f\^{}\{\textbackslash prime \textbackslash prime\}\textbackslash left(X\_\{s\^{}\{-\}\}\textbackslash right)\textbackslash left(\textbackslash Delta X\_\{s\}\textbackslash right)\^{}\{2\}\textbackslash right]

]



где $[X, X]=\left([X, X]_{t}, \mathscr{F}_{t}\right)$ - квадратич. вариация семимартингала $X$, то есть



\[

\begin{gathered}

{[X, X]_{t}=X_{0}^{2}+2 \int_{0}^{t} X_{s} d X_{s},} \\

{[X, X]_{t}^{c}=[X, X]_{t}-\sum_{0<s \leqslant t}\left(\Delta X_{s}\right)^{2}}

\end{gathered}

\]



\begin{itemize}

  \item непрерывная часть квадратич. вариации $[X, X]$, $\Delta X_{s}=X_{s}-X_{s-}$, а рассматриваемые интегралы понимаются как стохастич. интегралы по С.

\end{itemize}



Если $X$ есть С., то процесс $X^{(\leqslant 1)}=\left(X_{t}^{(\leqslant 1)}, \mathscr{F}_{t}\right)$ с



\[

X_{t}^{(\leqslant 1)}=X_{t}-\sum_{0<s \leqslant t} \Delta X_{s} I\left(\left|\Delta X_{s}\right|>1\right)

\]



имеет ограниченные скачки, $\left|\Delta X_{t}^{(\leqslant 1)}\right| \leqslant 1$, и в силу этого он допускает единственное представление вида



\[

X_{t}^{(\leqslant 1)}=X_{0}+B_{t}+M_{t},

\]



где $B=\left(B_{t}, \mathscr{T}_{t}\right)$ - предсказуемый процесс локально ограниченной вариации, а $M=\left(M_{t}, \mathscr{F}_{t}\right)$-локальный мартингал. Этот мартингал однозначным образом представим как $M=M^{c}+M^{d}$, где $M^{c}=\left(M_{t \mathscr{F}}^{c}\right)^{c}-$ непрерывный локальный мартингал (непрерывная мартингальная составляющая семимартингала $X$ ), а





\end{document}